\chapter{Configuration}
\label{ch:configuration}

Once the setup wizard is behind you, the day-one configuration work
consists of three blocks: defining your \emph{storage
locations}\index{storage locations}, populating the \emph{reference
data} (genres, contributor roles, volume states, location types),
and adjusting the system-wide \emph{settings} (overdue threshold,
auto-purge cadence, provider keys).

Everything in this chapter is configured from the
\texttt{/admin}\index{admin panel} panel — there is no separate
configuration file. The panel has five tabs:

\begin{itemize}
  \item \textbf{Health} — read-only status dashboard.
  \item \textbf{Users} — create / deactivate users, change roles.
  \item \textbf{Reference data} — the four taxonomy tables.
  \item \textbf{Trash} — view and restore soft-deleted items.
  \item \textbf{System} — global settings.
\end{itemize}

Only the \texttt{Admin} role sees the panel; \texttt{Librarian} and
\texttt{Anonymous} users get a 403 from every \texttt{/admin/*}
route.

\section{Storage locations and hierarchy}\label{sec:locations}

Storage locations are the physical spots where your volumes
live\index{volume}: a room, a shelf, a row, a numbered box, a
specific cell in a cabinet. mybibli organises them as a
\textbf{tree}\index{location tree} — every location has at most one
parent. The hierarchy is arbitrary: you decide both the levels and
the labels.

A workable household hierarchy might be three levels deep:

\begin{lstlisting}
Living room          (room)
 └── Bookshelf A      (shelf)
      └── Row 1       (row)
      └── Row 2       (row)
 └── Bookshelf B      (shelf)
      └── Row 1       (row)
Bedroom              (room)
 └── Nightstand      (shelf)
\end{lstlisting}

Three levels is enough to find any volume in seconds without forcing
the hierarchy to mirror every nook of your living space. A flat list
(one big ``shelf'' per physical shelf) works equally well for
collections under a few hundred volumes.

\subsection*{Editing the tree}

Open \texttt{/locations}. The tree is rendered as a nested list with
inline create / rename / move / delete affordances. Drag-and-drop is
\emph{not} supported (mybibli ships with a CSP-strict no-JS policy);
moves are done by editing a node and picking a new parent from a
dropdown.

\subsection*{Location node types}

Each location node carries a \textbf{node type}\index{node type} —
``room'', ``shelf'', ``row'', ``box'', anything you like. Node types
are managed under \texttt{/admin > Reference data > Location node
types}. They are pure labels; mybibli does not enforce that a
``shelf'' must sit under a ``room''.

\begin{tip}
Renaming a node type cascades to every location that uses it (the
type is stored by name in the locations table). Deleting a node type
is refused as long as one location still references it.
\end{tip}

\section{Labels: V-codes and L-codes}\label{sec:labels}

mybibli prints two kinds of barcode labels:

\begin{itemize}
  \item \textbf{V-codes}\index{V-code} — one per
        \emph{volume}\index{volume} (the physical book on a shelf).
        Format: \texttt{V} followed by a zero-padded sequence
        (e.g.\ \texttt{V0042}).
  \item \textbf{L-codes}\index{L-code} — one per
        \emph{storage location}. Format: \texttt{L} followed by a
        zero-padded sequence (e.g.\ \texttt{L0007}).
\end{itemize}

V-codes and L-codes are scanned with the same barcode scanner you
use for ISBNs at cataloging time. The scanner's prefix
(\texttt{V} or \texttt{L}) tells mybibli which lookup to perform.

\subsection*{Printing labels}

Open \texttt{/locations} or the relevant volume page and click
\textbf{Print label}. mybibli renders a printable HTML page with the
Code-128 barcode and human-readable text. Print on plain paper or
sticker stock — the encoding is forgiving enough that a low-cost
inkjet works for V-codes and L-codes alike.

\begin{tip}
Print V-code stickers \emph{before} you start cataloging in bulk. A
roll of pre-printed V0001 through V0500 means you can grab a sticker,
stick it on the back cover, and scan in one motion — no
context-switch between the cataloging UI and the printer.
\end{tip}

\section{Reference data}\label{sec:reference-data}

Four taxonomy tables drive the dropdowns you see while
cataloging. They live under \texttt{/admin > Reference data}:

\begin{description}
  \item[Genres]\index{genres} The genres assigned to titles. mybibli
        ships with a baseline FR/EN seed (Fiction, Non-fiction,
        Biography, \dots) — edit or extend to match your collection.
  \item[Contributor roles]\index{contributor roles} Author,
        Illustrator, Translator, Editor, Narrator, \dots Used by
        title detail pages to attribute the right people to the
        right slots.
  \item[Volume states]\index{volume state} The physical condition
        of an exemplaire — New, Good, Worn, Damaged, \dots Useful
        for loan returns (mark a volume Worn when the borrower
        damaged it).
  \item[Location node types]\index{node type} See
        section~\ref{sec:locations} above.
\end{description}

\subsection*{CRUD semantics}

All four tables share the same CRUD pattern:

\begin{itemize}
  \item \textbf{Create} adds a row. If the name collides with a
        soft-deleted row of the same table, the existing row is
        reactivated and the result reads \emph{Reactivated} instead
        of \emph{Created}.
  \item \textbf{Rename} updates the name. For location node types,
        the new name cascades to every location that uses the
        type.
  \item \textbf{Delete} performs a soft-delete \emph{only when the
        row is unused}. If any title still references the
        genre / role / state / type, the deletion is refused with
        a count of dependents.
\end{itemize}

\subsection*{Reference data is not translated}

Reference data values are stored verbatim — the FR seed creates a
genre called ``Roman'' and that is the literal text English-speaking
users will see if they switch the UI to English. Only the UI
chrome\index{i18n} (button labels, navigation, error messages) is
translated.

\section{System settings}\label{sec:system}

The \texttt{/admin > System} tab groups every global setting. The
panel is divided into three forms, each with its own save button:

\subsection*{Loans}

\begin{itemize}
  \item \textbf{Overdue threshold (days).} Defines how many days a
        loan can run before it surfaces as overdue on the dashboard.
        Default 30.
  \item \textbf{Session inactivity timeout (seconds).} Idle time
        after which an authenticated session is wiped. Default
        1800\,s (30\,minutes).
  \item \textbf{Auto-purge interval (seconds).} How often the
        background scheduler hard-deletes items older than 30\,days
        in the trash. Default 86400\,s (24\,hours).
\end{itemize}

\subsection*{Metadata providers}

One row per keyed provider (Google Books, OMDb, TMDb). Each row has
a key input and a \emph{Clear} button. Saving a key takes effect on
the next metadata fetch — no restart, no re-deploy.

\begin{note}
If the same provider key is set in \texttt{.env} \emph{and} cleared
from the UI, the next reboot re-migrates the value from
\texttt{.env}. The operator's deployment-time intent wins on boot.
To make a Clear durable across reboots, also remove the variable
from \texttt{.env} (or \texttt{docker-compose.yml}).
\end{note}

\subsection*{Language}

Sets the default UI language for \emph{anonymous} visitors.
Authenticated users override this with the language toggle in the
navigation bar — the choice is stored on their user row.

\section{Users}\label{sec:users}

Open \texttt{/admin > Users} to manage user accounts. Each row shows
the username, role, active flag, and last-seen timestamp. The
available actions are:

\begin{itemize}
  \item \textbf{Create user} (top of the page) — username, password,
        role.
  \item \textbf{Edit} — change role or rename. Cannot change another
        user's password from this view; password resets go through
        the user's own \texttt{Account} page.
  \item \textbf{Deactivate} — soft-delete the user. Their sessions
        are wiped immediately; they cannot log back in.
  \item \textbf{Reactivate} — undo a deactivation (from the
        \texttt{Trash} tab).
\end{itemize}

Two safety guards apply:

\begin{itemize}
  \item \textbf{Self-deactivate is blocked.} You cannot deactivate
        the account you are logged in as.
  \item \textbf{Last-active-admin guard.} The system always keeps at
        least one active admin. Deactivating the last one is
        refused with a 409 conflict and a clear error.
\end{itemize}

The three roles\index{roles} (Anonymous, Librarian, Admin) are
detailed in chapter~\ref{ch:roles}.
